% Part: model-theory
% Chapter: interpolation
% Section: introduction

\documentclass[../../../include/open-logic-section]{subfiles}

\begin{document}

\olfileid{mod}{int}{int}

\olsection{引言}

插值定理如下：假设
$\Entails !A \lif !B$，则存在语句 $!C$，使得
$\Entails !A \lif !C$ 且 $\Entails !C \lif !B$。
此外，$!C$ 中出现的所有常项符号、函数符号和谓词符号（除 $\eq$ 外），
既在 $!A$ 中出现，也在 $!B$ 中出现。
语句 $!C$ 称为 $!A$ 与 $!B$ 的一个\emph{插值式}。

插值定理本身就颇具意义，但其主要重要性在于，
它可用来证明理论中可定义性的相关结论，
以及两个一致理论合并后仍为一致理论的条件。
前一结论称为 Beth 可定义性定理；后一结论称为 Robinson 联合一致性定理。

\end{document}